\scanpage{4}
% 1–3
est la condition

e) Le composé de deux \inserted{$A_*$-}morphismes $X_m-X_n-X_p$ est un $A_*$-morphisme.

% 3, left
% slanted
\margin{Moyennant a) à e) $\mathrm{Hom}_A(X_m,X_n)\ \unsure{\simeq|\approx}\ \mathrm{Hom}_{A_0\text{-alg}}(A_n,A_m)$ \unsure{?}}

% 4–8
\uline{NB} Moyennant \inserted{b) et} e), la condition a) équivaut au fait que
\[
\left\{\begin{aligned}
 a_1)&\quad 1\in A_0\\
 a_2)&\quad \text{le morphisme ``différence}\quad X_1\times X_1-X_1\text{ est } A_*\text{-morph.}\\
 a_3)&\quad \text{--------------- ``produit''}\quad X_1\times X_1-X_1\text{ est } A_*\text{-morph.}
\end{aligned}\right.
\]

% 9–10
\uline{NB} Les $A_n=k[t_1,\ldots,t_n]$ ou $A_n=k[t_1,\ldots,t_n]_{\mathfrak m_n}$ satisfont e)

% 11–15
L'analogue de c) est la condition

f) Si $u:X_n\longrightarrow X_n$ est un $A$-morphisme et un isomorphisme, \deleted{alors} i.e. le déterminant \deleted{\unsure{fct|fon}} de la matrice jacobienne de $u$ est inversible, alors $u^{-1}$ est un $A$-morphisme

% 16–17
Cette condition implique \deleted{(\unsure{et est|et est ce} équivalente à ?)}

% 16, left
% slanted
\margin{f') implique-t-il f (moyennant a) à e)) ? il ne semble pas}

% 18–20
f') Les $A_n$ sont des anneaux henséliens

[En fait, si f) est satisfait pour un $n$, (f') est satisfait pour cet $n$, i.e. $A_n$, est hensélien.]

% 20–24
\deleted{D'ailleurs si $u:\inserted{X_n\xrightarrow{\sim}X_n}$ est donné par $f_1,\ldots,f_n\in B_n$, alors $u^{-1}$ est donné par $g_1,\ldots,g_n$ appartenant à la clôture hensélienne de $(k[f_1,\ldots,f_n]_{\mathfrak n}$, où $\mathfrak n=k[f_1,\ldots,f_n]$.}

% 22, margin
% slanted
\margin{\deleted{\unsure{vérifier}}}

% 25–28
\uline{NB} Le plus petit $A_*$ satisfaisant a) à d) et f') est formé des
\[
 B_n=\mathbb P\{t_1,\ldots,t_n\}\qquad
 \left(\begin{aligned}&\text{clôture hensélienne de}\\[-2pt]&\mathbb P[t_1,\ldots,t_n]_{\mathfrak m_n}\end{aligned}\right)
\]
et il satisfait $=$ e) et f).

% 29
\deleted{g) Soit $\xi$ un champ de vecteurs \unsure{for|fon}}

% 30–33
Sous a) b) c) d) e) on définit la catégorie des $A_*$-schémas formels\inserted{et des Modules dessus, fibrés de t.f. !}, et le calcul différentiel dessus (opérateurs différentiels, parties principales, \inserted{fibré tangent} etc ....)

% 31, left
% slanted
\margin{\unsure{(des graphismes)|des graphismes)|6 des propriétés} \unsure{Th.|Obs.} des produits fibrés transversaux etc ....}
% archive sheet number: 3
